Un capteur solaire plan d'aire \qty{\v{S}}{m^2} reçoit perpendiculairement un rayonnement solaire de \qty{\v{I}}{W/m^2}. Sa surface sélective absorbe une fraction \num{\v{ab}} du rayonnement solaire, mais n'a qu'une émissivité de \num{\v{em}} dans l'infrarouge. L'environnement est à \qty{\v{Te}}{\degreeCelsius}. On suppose qu'aucune énergie n'est retirée du capteur et que seule sa face avant rayonne, sans convection.

\begin{enumerate}[label=\alph*)]
    \item Déterminez la puissance solaire absorbée par le capteur.
    \item Déterminez la température d'équilibre du capteur.
    \item Pourquoi les fabricants cherchent-ils une absorptivité élevée dans le visible et une émissivité faible dans l'infrarouge ?
\end{enumerate}
